\documentclass[11pt,a4paper]{article}
\usepackage{templates/modeling-notes}
\hypersetup{pdftitle={MATLAB review 1 -- Arrays and linear systems}}
\title{\textbf{MATLAB Review}\\[2mm]\large 1: Arrays, Indexing, and Linear Systems}
\author{Personal study notes}
\date{}
\begin{document}
\maketitle
\noindent\textbf{Purpose:} Revisit the matrix operations needed for Lecture 2.
\tableofcontents

\section{Rows, columns, and dimensions}\label{sec:arrays}
A decision vector is usually a column. Commas or spaces separate entries within
a row; semicolons separate rows. A semicolon after a statement suppresses its output.
\begin{lstlisting}
x = [2; 6];                 % 2-by-1 column
c = [3; 2];                 % 2-by-1 column
M = [1, 2, 3; 4, 5, 6];    % 2-by-3 matrix
size(M)                     % Expected: [2 3]
numel(M)                    % Expected: 6
revenue = c.' * x;           % Expected: 18
\end{lstlisting}
The revenue product has dimensions $(1\times2)(2\times1)=1\times1$.
\texttt{.'} transposes without conjugating; \texttt{'} also conjugates complex
entries. They agree for real data. Prefer \texttt{size} for matrix dimensions:
\texttt{length} gives the largest dimension and does not count all elements.

\section{Indexing and selections}\label{sec:indexing}
Indices start at 1. The first index selects the row; the second selects the column.
A colon selects all entries along that dimension. Linear indexing reads columns
from top to bottom, then proceeds to the next column.
\begin{lstlisting}
M(2, 3)                    % 6
M(:, 2)                    % [2; 5]
M(1, :)                    % [1 2 3]
M(:)                       % [1; 4; 2; 5; 3; 6]
M(end, :)                  % Last row: [4 5 6]
M(M > 3)                   % [4; 5; 6]
\end{lstlisting}
Logical indexing selects entries satisfying a condition. For a collection of
candidate decision vectors stored as columns, select whole columns with a logical
mask rather than extracting individual entries.

\section{Matrix operations versus element-wise operations}\label{sec:operations}
\begin{lstlisting}
M = [1, 2; 3, 4];
M * M                      % [7 10; 15 22]
M .* M                     % [1 4; 9 16]
M ^ 2                      % Same as M*M
M .^ 2                     % Same as M.*M
\end{lstlisting}
For matrix multiplication,
$(M^2)_{ij}=\sum_k M_{ik}M_{kj}$; for element-wise squaring,
$(M\mathbin{.\!\wedge}2)_{ij}=M_{ij}^2$. The dot changes the mathematical operation.
For sampled function values, typically use \texttt{.*}, \texttt{./}, and
\texttt{.\^}; for a linear model $Ax$, use matrix multiplication.

\begin{remark}[Implicit expansion can hide shape mistakes]
Adding a $2\times1$ column to a $1\times2$ row produces a $2\times2$ array by
implicit expansion. It need not raise a dimension error. Inspect dimensions when
an objective or feasibility calculation unexpectedly returns a matrix.
\end{remark}

\section{Solve a small system and check the model}\label{sec:solve}
The first two tight production constraints from Lecture 2 give
\[
2x_1+x_2=10,\qquad 3x_1+3x_2=24.
\]
Dividing the second equation by 3 gives $x_1+x_2=8$. Subtraction yields $x_1=2$,
then $x_2=6$. MATLAB expresses this system as:
\begin{lstlisting}
B = [2, 1; 3, 3];
b = [10; 24];
x = B \ b;                 % Expected: [2; 6]
residual = norm(B*x - b, Inf);
assert(residual <= 1e-10);
\end{lstlisting}
Backslash solves the linear system. For this nonsingular square $B$, the solution
is unique; avoid calculating \texttt{inv(B)*b}. For general matrices, inspect rank,
conditioning, warnings, and residuals: a small residual alone does not prove a
small error in the solution.

Check \emph{all} constraints of the original model as well:
\begin{lstlisting}
A = [2, 1; 3, 3; 1.5, 0];
capacity = [10; 24; 6];
slack = capacity - A*x;     % Expected: [0; 0; 3]
assert(all(x >= -1e-10));
assert(all(slack >= -1e-10));
\end{lstlisting}
These verify feasibility. They do not search for an LP optimum; Lecture 2's
hand certificate proves the objective cannot exceed 18 PLN.
For these small examples $10^{-10}$ is an absolute tolerance. Real problems
require tolerances chosen with their units, scaling, and numerical conditioning.

\section{Function handles, a loop, and a plot}\label{sec:functions}
Consider the quadratic
$f(t)=t^2-4t+5=(t-2)^2+1$. Its global minimum is $f(2)=1$.
\begin{lstlisting}
f = @(t) t.^2 - 4*t + 5;
t = linspace(0, 4, 101);
values = f(t);
loopValues = zeros(size(t));
for k = 1:numel(t)
    loopValues(k) = f(t(k));
end
assert(max(abs(values-loopValues)) <= 1e-10);
\end{lstlisting}
The function handle can be called on a scalar or the entire sample vector because
the square is element-wise. Preallocate the loop output instead of repeatedly
growing an array. Agreement between the loop and vectorized evaluation checks
implementation consistency; the completing-the-square argument establishes the minimum.

\begin{lstlisting}
fig = figure('Visible', 'off');
plot(t, values, 'LineWidth', 1.5);
xlabel('t'); ylabel('f(t)');
title('f(t) = (t - 2)^2 + 1'); grid on;
exportgraphics(fig, 'quadratic-review.png', 'Resolution', 150);
close(fig);
\end{lstlisting}
The accompanying script saves its figure to the \path{matlab/output/}
directory using the script location to find the repository root.
A sampled graph illustrates the function; it does not replace a proof or guarantee
that a minimum lies on the sampled grid.

\section{Run the checks and practise}\label{sec:practice}
From the repository root, run:
\begin{lstlisting}
run('matlab/review/arrays_and_linear_systems.m')
\end{lstlisting}
The script contains assertions for the hand results above. Before running, predict the results, then explain any difference.
\begin{enumerate}[leftmargin=1.8em]
\item Write a column vector with entries 1 through 5 and select its last two entries.
\item Compute both the matrix square and element-wise square of a $2\times2$ matrix.
\item Replace the production candidate by $(4,4)^T$. Which resource constraint fails?
\item Change the quadratic to $t^2-6t+10$. Derive its minimum by hand before plotting.
\end{enumerate}

\section{Sources and editorial remarks}\label{sec:sources}
\begin{itemize}[leftmargin=1.6em]
\item \href{https://www.mathworks.com/help/matlab/learn_matlab/matrices-and-arrays.html}{MathWorks: Matrices and Arrays}:
array construction, dimensions, and transpose.
\item \href{https://www.mathworks.com/help/matlab/math/array-indexing.html}{MathWorks: Array Indexing}:
subscripts, linear indices, and logical masks.
\item \href{https://www.mathworks.com/help/matlab/matlab_prog/array-vs-matrix-operations.html}{MathWorks: Array vs. Matrix Operations}:
operator distinctions and compatible array sizes.
\item \href{https://www.mathworks.com/help/matlab/ref/double.mldivide.html}{MathWorks: mldivide}:
linear system solution and numerical warnings.
\item \href{https://www.mathworks.com/help/matlab/ref/exportgraphics.html}{MathWorks: exportgraphics}:
graphics export.
\end{itemize}
Polish terms are maintained in \texttt{docs/glossary.md}: \emph{wektor wierszowy},
\emph{wektor kolumnowy}, \emph{indeksowanie}, \emph{operacje elementowe},
\emph{uchwyt do funkcji}, and \emph{reszta uk\l{}adu r\'owna\'n}.
\paragraph{Provenance and verification.}
\textbf{Sources:} Official MathWorks documentation listed at the end;
Lecture 1's historical refresher scope motivates this topic.\\
\textbf{Execution status:} Verified on 2026-10-04 in MATLAB
26.2.0.3386108 (R2026b), based on output supplied by the user from a normal terminal.
All numerical assertions passed and the plot was exported.
Hand results agree with the script assertions.

The practice prompts in Section~\ref{sec:practice} are original examples.
Full terminal output is recorded in \texttt{docs/matlab-verification.md}.
\end{document}
